\section{Three connected research questions}
Small farmers may pool products and sell through an intermediary who reaches households in underserved markets. The arrangement reduces search costs, but the intermediary's margin rule divides surplus and shapes effort and product selection. Coase explains why transaction costs shape organization~\cite{coase1937,nobelcoase1991}. Morgenstern et al. show that resale-platform margin competition can reduce product-search effort and that centralized margin decisions shift competition toward product selection~\cite{morgenstern2025}. I adapt this mechanism to farmer sales; their paper does not establish rural-consumption-inequality effects.

\textbf{Q1---Economics:} How do fixed and intermediary-chosen margins affect farmer net earnings and household product access when effort and coordination costs matter?\\
\textbf{Q2---Computation:} Under which demand, cost, competition, and trust conditions does each rule improve participation, completed purchases, and surplus?\\
\textbf{Q3---Behavior:} Does reputation information raise trust and willingness to pay enough to change which rule performs better?\\
\textbf{Integrated question:} Which rule improves farmer earnings and real household access after effort, product selection, coordination cost, and trust are included? Figure~\ref{fig:teaser} displays the causal chain.

\section{Answer to the economics question}
The players are a platform, a pooled farmer group, an intermediary, and households. The platform selects $R\in\{F,C\}$; farmers choose participation and pay coordination cost $K$; under $F$ the platform fixes margin $m$, while under $C$ the intermediary chooses $m$ and effort $e$. Let $w$ be the farmer payment, $c$ unit cost, $p=w+m$, and $q(m,e,r)$ completed purchases, with $r$ a reputation signal. Payoffs are
\[\pi_F=(w-c)q-K,\qquad \pi_I=mq-g(e).\]
I solve this sequential game by backward induction. The testable prediction is a tradeoff: a moderate fixed margin can lower retail price and raise purchases, while a chosen margin can increase intermediary profit and effort; very low margins may cause exit. This is a hypothesis, not an assumed benefit of collective selling.

\section{Answer to the computational question}
The notebook implements a computational grid search. It searches feasible $(m,e)$ choices, applies farmer participation, and reports price, completed quantity, farmer earnings, intermediary profit, and surplus. With $q=12e\max\{0,4-m\}$, $w=6$, $c=4$, $K=12$, $g(e)=30e^2$, and fixed $m=1$, the synthetic baseline gives fixed $(e,q,\pi_F)=(0.60,21.60,31.20)$. The chosen-margin grid gives $(m,e,q,\pi_F)=(2.00,0.80,19.20,26.40)$ and intermediary profit $19.20$ versus $10.80$ under the fixed rule. These are synthetic outputs, not field findings. Sensitivity analysis varies coordination cost, demand, trust, effort cost, competition, and heterogeneous products; variety is a planned extension.

\clearpage\balance
\section{Answer to the behavioral question}
Households may not convert listed availability into purchases if they distrust an unfamiliar seller. Building on reputation experiments~\cite{resnick2006} and bounded rationality~\cite{simon1955}, I will estimate purchase probability at a given posted price under sparse versus verified reputation information. Purchase, trust, and willingness-to-pay measures distinguish a trust mechanism from a price-only explanation. The estimated response replaces $r$ in $q(m,e,r)$ and the simulation is rerun. A small online study cannot identify long-run farmer income or population welfare; a result showing no reputation effect or a consistent advantage for chosen margins would revise my hypothesis.

\section{Advanced development and future directions}
This proposal extends Coase's transaction-cost question~\cite{coase1937,nobelcoase1991} to digital intermediation between farmers and underserved households. Newell and Simon's Turing-recognized work on computational decision processes~\cite{turing1975} motivates an inspectable model of bounded choices. Today's platforms can jointly set margins, rank products, and display reputation. Current resale-platform theory does not determine whether these rules raise farmer earnings after coordination costs or turn listed products into trusted completed purchases.
\begin{table}[h]\caption{Roadmap from foundation to contribution.}\begin{tabularx}{\columnwidth}{@{}p{1.25cm}Y@{}}\toprule\textbf{Stage}&\textbf{Milestone}\\\midrule Foundation&Coase on organization; Newell and Simon on computational decisions.\\ PS1&Fixed/chosen margin baseline with coordination cost.\\ Next&Trust evidence, product heterogeneity, competition, and field validation.\\ 2056&Adaptive, auditable rules that improve access and producer welfare.\\\bottomrule\end{tabularx}\end{table}
\noindent\textbf{Expected 2056 Nobel Prize and Turing Award contribution statement.} By 2056, I aspire to develop a general theory and open computational system for designing intermediated markets under transaction costs and bounded trust. I will connect economic incentives, auditable simulation, and behavioral evidence to identify rules that improve access without weakening producer participation.
\subsection*{Open Science Statement}\textbf{GitHub:} \GitHubLink\\\textbf{Colab:} \ColabLink\\\textbf{Submitted commit:} \SubmittedCommit\\The repository contains the model, requirements, notebook, tests, synthetic inputs, and saved outputs. A clean run reproduces the baseline and sensitivity files.
\subsection*{Statement of Contribution to the UN's SDGs}\textbf{Hugging Face:} \DemoLink\\The planned Margin Rule Lab will support \textbf{SDG 4---Quality Education}~\cite{sdg4} by letting learners vary margins, coordination costs, effort costs, and reputation and observe synthetic access and earnings outcomes. Educational benefit will be checked with a short pre/post question.
